% =============================================================================
% ABSTRACT PAGE
% =============================================================================
\clearpage
\addcontentsline{toc}{chapter}{Abstract}
\begin{center}
{\large\bfseries Abstract of the Dissertation}

\vspace{0.3in}

{\bfseries Three Essays on the Econometric Analysis\\
of Emerging Issues in Electricity Markets}

\vspace{0.2in}

by

\vspace{0.2in}

{\bfseries Dana Golden}

\vspace{0.2in}

{\bfseries Doctor of Philosophy}

\vspace{0.1in}

in

\vspace{0.1in}

{\bfseries Economics}

\vspace{0.1in}

Stony Brook University

\vspace{0.1in}

2026
\end{center}

\vspace{0.3in}

\doublespacing
This dissertation comprises three essays. Each essay examines an emerging issue in the electricity market and impacting the energy transition using econometric methods. The dissertation maps the challenges facing the market up and down the supply chain and provides counterfactual insights for incorporation by policy makers.

The first essay estimates how transmission expansion affects generator investment in electricity markets. Renewable resources are essential for decarbonization, but their unequal geographic distribution limits adoption. I estimate a dynamic model of generator behavior combining a short-run dispatch problem with transmission constraints and a long-run dynamic investment game, using an originally constructed dataset combining EIA, ISO, and proprietary data from 2018–2023. I find that upgrading inter-zonal transmission to modern HVDC standards would reduce average wholesale prices by 6\%, generating approximately \$4 billion in annual welfare gains, with heterogeneous regional effects. Major transmission projects such as the Grain Belt Express increase solar and wind adoption in the Eastern Interconnection by over 10\% by 2050. I use my estimates to evaluate the Big Wires Act, the Inflation Reduction Act, and the Bipartisan Infrastructure Law.


This second essay studies how learning-by-doing shapes cost dynamics, technological adoption, and the effectiveness of industrial policies in the solar panel industry. We develop and estimate a dynamic model of firm entry, exit, production, and technology choice using firm-level data on Chinese manufacturers and global data from 2000 to 2013. The model allows production experience to reduce future costs, generating a dynamic feedback loop between output and productivity.
We find that learning-by-doing is a key driver of cost reductions, particularly for crystalline silicon (c-Si) technology: doubling cumulative production reduces marginal costs significantly. Learning also amplifies the effects of industrial policies, leading to larger long-run impacts on output and prices. It further shapes technological competition, contributing to the dominance of c-Si technology. Finally, earlier policy implementation yields substantially larger effects by accelerating movement along the experience curve.

The third essay measures the grid-level impacts of artificial intelligence. We develop a game-theoretic model identifying a timing wedge, the mismatch between data center consumption and credited renewable generation, as the central mechanism through which AI demand degrades reliability, raises prices, and increases emissions even under full REC coverage. AI demand significantly increases fossil generation, wholesale prices (up to 25\% in treated PJM zones), and outage frequency (0.5–1 additional outages per year) near data centers, with impacts scaling in model size. Data centers with on-site generation exhibit a sign reversal in power-quality effects, consistent with the model's prediction that behind-the-meter capacity absorbs demand spikes.

\clearpage
